\chapter{Conclusion and Future Work}

This chapter closes the report. It first restates what was actually built and
maps it back to the objectives set out in Chapter~1, then records the lessons
that came out of building it, and finally lays out the work that would take the
system from a working platform to something a seed cooperative could run in
production for a season.

\section{Summary of Achievements}

The project set out to build a seed-quality intelligence platform: a user
photographs a batch of seeds, and the system detects each individual seed and
classifies its quality, then reports aggregate metrics such as seed count,
good-rate, and a per-crop breakdown. That platform exists and runs. What follows
is an account of the pieces that were delivered, organised around the objectives
from Chapter~1.

\subsection{An end-to-end, traceable inference platform}

The core of the product is a two-stage vision pipeline. A detection model scans
the whole image and returns a bounding box, a confidence, and a crop class for
every seed it finds; then each detected box is cut from the source image and fed
to a binary classifier that returns \texttt{good} or \texttt{bad} with a
confidence \cite{kamilaris2018deepag}. The detector is a torchvision Faster
R-CNN \cite{fasterrcnn2015} with a ResNet-50 backbone and a Feature Pyramid
Network \cite{resnet2015, lin2017fpn}, carrying a three-class head over
\texttt{background}, \texttt{coffee}, and \texttt{maize}. The classifiers are
ResNet-18 backbones with a Convolutional Block Attention Module inserted after
the final residual stage \cite{cbam2018}, one classifier trained per crop type.
Splitting the problem this way means a detector and a classifier can be
versioned, swapped, and A/B-tested independently, because detection and
classification are recorded as separate \texttt{inferences} rows over the same
image.

Traceability was a hard requirement rather than a nice-to-have, and it is
enforced at the schema level. Every \texttt{seed\_detection} row carries a
foreign key to the \texttt{inference} that produced it, and every
\texttt{inference} carries a foreign key to the exact \texttt{model\_artifact}
that ran. Any single seed label in the system can be traced back to the model
version and weights that generated it. This chain is what makes the rest of the
ML platform trustworthy: an A/B result or an offline evaluation score means
nothing if you cannot say which weights produced which label.

The request path that ties these together is the unified \texttt{POST /analyze}
endpoint. An upload lands in MinIO and Postgres, a Celery task picks it up
\cite{celeryinplay}, the traffic router selects a model for the segment, the
detector runs, each crop is classified, and the results are persisted with the
traceability chain intact. Batch state moves through a state machine
(\texttt{pending} $\rightarrow$ \texttt{running} $\rightarrow$ \texttt{succeeded}/\texttt{partial}/\texttt{failed}),
and clients poll \texttt{GET /batches/\{id\}} for progress. The failure handling
is deliberate: if classification crashes after detection has already persisted,
the batch degrades to \texttt{partial} rather than discarding good detection
results, and if no model is promoted the user gets a readable message instead of
a blank error.

\subsection{A model registry with A/B routing and offline evaluation}

Weights are uploaded to MinIO and recorded in the \texttt{model\_artifacts}
table, each row carrying its kind, backend, builder key, a per-model config
blob, and a lifecycle status that moves
\texttt{registered} $\rightarrow$ \texttt{staging} $\rightarrow$ \texttt{production} $\rightarrow$ \texttt{archived}.
Only staging and production models can serve traffic. The traffic router decides
which model runs for a given segment: when weighted splits exist in
\texttt{traffic\_splits} it routes by a sticky per-user bucket so the same user
always lands on the same model during an experiment, and when no split exists it
falls back to the production model for that crop type and then to a global
production model. That fallback is what lets the mobile point-and-shoot flow,
which sends no crop type, still resolve to a model.

Around the registry sits the experiments subsystem. A labelled dataset stored in
MinIO and Postgres is the ground truth; an experiment runs a chosen model over
that dataset offline, computes classification metrics including the confusion
matrix, logs parameters, metrics, and artifacts to MLflow \cite{mlflow2018}, and
writes a Markdown report. The model manager keeps a per-process LRU cache of
loaded modules, serialises concurrent loads of the same model through a lock so
two workers never duplicate a GPU load, and hot-reloads a model when its artifact
timestamp advances. A developer can register new weights, evaluate them offline,
promote them, and route a slice of live traffic to them without a restart.

\subsection{Three clients, bilingual and right-to-left}

The backend serves one API to three clients. The first client is the HTTP API
itself, a versioned, layered FastAPI service \cite{fastapi, fielding2000rest}
with JWT and OAuth2 authentication \cite{rfc7519jwt, rfc6749oauth}, role-based
access control, per-route rate limiting, and RFC~9457 problem-details errors
\cite{rfc9457problem}. The second is a React web application
\cite{reactindepth} covering the farmer journey of sign in, photograph, report,
and history, plus the developer-facing model and analytics pages. The third is an
Expo / React Native mobile app \cite{reactnativeexpo} with a realtime camera
capture flow. Both user-facing clients ship in Arabic and English with full
right-to-left layout, which matters for the intended Egyptian farming audience
and is not something that can be retrofitted cheaply.

\subsection{Observability and operations}

The platform is instrumented rather than opaque. It exposes Prometheus metrics, a
Grafana dashboard set, OpenTelemetry tracing \cite{opentelemetry2022}, Sentry
error reporting, structured logging, and \texttt{/healthz} and \texttt{/readyz}
probes. Configuration comes from one central \texttt{Settings} object sourced
from environment variables or file-based secrets, never read piecemeal from the
process environment. The whole stack ships as a multi-stage Docker build
\cite{merkel2014docker} with a development compose file and a production overlay.
ClickHouse \cite{clickhouseolap} backs the analytics pages through a star schema
so heavy analytical queries stay off the OLTP Postgres \cite{postgresinprod}.

\subsection{MultiSeedGen: closing the labelling gap}

The companion tool, MultiSeedGen, addresses a problem that blocks the detector
more than any architectural choice does. Training an object detector needs images
of \emph{many} seeds with a bounding box drawn around each one, and labelling
those boxes by hand is slow and error-prone. What is comparatively easy to obtain
is photographs of single seeds. MultiSeedGen turns the cheap input into the
expensive output: it segments individual seeds out of source photos using a
cascade of classical masks and learned matting via rembg and an optional Segment
Anything backend \cite{qin2020u2net, kirillov2023sam}, then composites many of
those cutouts onto realistic backgrounds and exports a detection dataset in YOLO,
YOLO-seg, and COCO formats with the bounding boxes generated for free as a
by-product of placement \cite{dwibedi2017cutpaste, lin2014coco, yolov8}.

Two design decisions make the output usable for training rather than a toy.
First, the compositor applies domain-matched lighting and camera degradation, a
deliberate move toward closing the gap between synthetic and real images through
domain randomisation and augmentation \cite{tobin2017domainrand,
shorten2019augmentation, buslaev2020albumentations}. Second, output is
byte-reproducible for a fixed configuration and seed, so a dataset is defined
entirely by its recipe; an experiment can be rerun exactly, which is the same
reproducibility discipline the inference platform enforces through its
traceability chain. The tool ships with a CLI and a FastAPI web UI over the same
pipeline, and writes provenance and QA reports alongside every run.

Taken together, the two systems form a loop on paper: MultiSeedGen produces the
labelled detection data the platform's detector is trained on, and the platform
produces the predictions that, in the future, can tell MultiSeedGen which kinds
of scenes it is generating too few of. Section~\ref{sec:ch7-future} returns to
closing that loop in practice.

\section{Lessons Learned}

The technical lessons that came out of this project were mostly about discipline
paying off slowly, and shortcuts charging interest later.

The strict layered architecture earned its keep. Holding the line on
\texttt{routers} $\rightarrow$ \texttt{services} $\rightarrow$ \texttt{repositories} $\rightarrow$ ORM
\cite{cleanarchitecture2017}, with Pydantic schemas at every boundary and the ORM
never leaking past the service layer, sounded like ceremony at the start. Its
value showed up when the ML platform grew. Because every backend satisfies the
same duck-typed protocol and returns framework-free DTOs, adding the YOLO and
Roboflow backends alongside the default local-torch one was a matter of writing a
class and registering it in a factory \cite{designpatterns1994},
with no churn in the services that call them. The one place that skipped the
pattern, the traffic router querying \texttt{traffic\_splits} ad-hoc instead of
through a repository, is also the one place that now reads as a wart, which is its
own kind of confirmation.

The single \texttt{Settings} object was a similar bet that paid off. Routing all
configuration through one validated object, sourced from environment variables or
file secrets, meant the production compose overlay could swap in file-based
secrets without touching application code, and it removed the class of bug where
two modules read the same setting differently. It is a small rule, applied
everywhere, and the payoff is that there is exactly one place to look.

Keeping torch out of the API image turned out to matter more than expected. The
heavy imports for torch and ultralytics live inside method bodies in the
inference backends, so the API process never imports torch; only the inference
worker does. The benefit is a small API container that starts fast and a clean
separation between the web tier and the GPU tier. The cost is that this invariant
is easy to break by accident with a careless top-level import, which is why one
of the planned CI checks is an explicit assertion that the API image contains no
torch \cite{sqlalchemyguide}.

The most expensive lesson was about tests. The suite is not hermetic: the
app-booting integration tests fail outside Docker because the rate limiter dials
a real Redis at the compose hostname, and the fakeredis fixture does not
intercept that client \cite{redisinaction}. Without continuous integration to
catch it, the suite rotted quietly. A fresh run surfaced two stale assertions
that referenced renamed internals, and 13 of the integration tests failed for
this Redis dependency rather than any product defect. The honest position is that
the committed coverage number is a partial artifact and cannot be trusted until
the suite runs the same way everywhere. Non-hermetic tests do not announce
themselves; they pass on the author's machine and quietly stop meaning anything.

The synthetic-to-real domain gap is real and was treated as such. Generating
plausible-looking multi-seed images is the easy part. Generating images close
enough to real field photographs that a detector trained on them transfers
cleanly is the hard part, which is exactly why MultiSeedGen invests in
camera-matched degradation and domain randomisation rather than clean composites.
Synthetic data narrows the labelling bottleneck; it does not erase the
distribution shift, and the right way to measure whether it has helped is to
evaluate the detector against real labelled scans, not against more synthetic
data.

Finally, coordinating a backend, two frontends, and a separate data tool was a
project-management lesson as much as a technical one. Two of the three planned
clients, two crop classifiers, a detector, and a generator each move at their own
pace, and the only thing that kept them coherent was treating the API contract
and the data formats as the fixed interfaces and letting everything else vary
behind them. The places where that discipline slipped, such as documentation
pointing at files that had moved and a script referenced in the README that was
never written, are precisely the places where the contract was implicit rather
than enforced.

\section{Future Enhancements}
\label{sec:ch7-future}

The work below is ordered roughly by how much it unblocks. The first item is a
precondition for trusting almost everything else.

\paragraph{Make the test suite hermetic and publish a real coverage number.}
The integration tests must run identically on a developer laptop and in CI. That
means making the application client use fakeredis end to end, or adding a Redis
service to the test compose file, so the rate limiter has something to dial
without a live cluster. A dependency lockfile should be added in the same pass so
installs are reproducible and library drift stops silently breaking the build.
Only after that is continuous integration worth standing up, and only then is a
coverage number meaningful enough to put in a report. The worst-covered areas are
already known: the worker tasks and the ML backends are barely exercised, and
MinIO storage and the OAuth callback are mocked everywhere.

\paragraph{Presigned and resumable uploads, and batch cancellation.}
Inference uploads are multipart-only today. Large field photographs and slow
rural connections argue for presigned uploads that go straight to MinIO, with
resumable transfers so a dropped connection does not restart a multi-megabyte
upload from zero. The service and schema for presigned upload already exist; what
is missing is the route. Batch cancellation belongs in the same effort: a user
who realises they uploaded the wrong batch should be able to stop a running scan
rather than wait for it to finish.

\paragraph{True change-data-capture to the warehouse.}
The ClickHouse warehouse is populated by an application-level dual-write: when a
fact is written to Postgres, a Celery task also writes it to the star schema.
This works but couples analytics correctness to application code and can drift if
a write path is added without a matching dual-write. Postgres is already
configured for logical replication; the enhancement is to consume that
replication stream and let change-data-capture populate the warehouse, so
analytics derive from the source of truth instead of a parallel write that has to
be kept honest by hand.

\paragraph{More crop types.}
The platform supports coffee and maize today, and the architecture was built to
extend. A new crop is a new detector head class, or a retrained detector, plus a
new per-crop classifier registered through a builder and the model registry. The
traffic router already routes per crop type, so a new classifier can be promoted
and A/B-tested in isolation without disturbing the existing crops. On the data
side, MultiSeedGen extends the same way: a new crop is a new set of source seed
photos and a config recipe, with no pipeline changes.

\paragraph{On-device inference for the mobile app.}
The mobile app currently captures images and sends them to the backend for
inference. For field use with poor connectivity, running a lightweight detector
and classifier on the device would let a farmer get a result with no round trip.
This is a natural fit for a quantised YOLO detector \cite{redmon2016yolo,
yolov8} exported to a mobile runtime, with the server remaining the authority for
history, A/B routing, and the traceability record once connectivity returns.

\paragraph{An active-learning loop.}
The most valuable enhancement closes the loop between the two systems. Real scans
that the platform classifies with low confidence, or that a user corrects, are
exactly the cases the current models handle worst, and they describe the scenes
the synthetic generator is underproducing. The loop is: collect those hard real
scans through the existing traceability chain, feed their characteristics back
into MultiSeedGen to regenerate training data weighted toward the failure modes,
retrain the detector and the affected classifier on the augmented set, evaluate
the candidate offline against a held-out set of real labelled scans, and promote
it through the registry's existing staging-to-production lifecycle behind a
traffic split for safety. Each turn of the loop targets generation at the
platform's measured weaknesses instead of producing more of what it already
handles well. The registry, the offline-evaluation runner, the traffic router,
and MultiSeedGen's reproducible recipes are all already in place; what remains is
the orchestration that connects them.
